\chapter{The Habitable Zone and the HZIED Test}
\label{chap:exoplanet_characterization}

This chapter sets up the physics the rest of the thesis will use.
\Cref{sec:habitable_zone} defines the classical habitable zone and
locates its inner edge at the runaway greenhouse threshold.
\Cref{sec:hzied_framework} defines the \HZIED test itself --- the
specific hypothesis that \textcite{Schlecker2024} proposed and this
thesis extends.  The observational technique (transit photometry)
and the empirical demographic backdrop that the test is applied to
are the subject of \Cref{chap:kepler_tess}.


% =======================================================================
\section{Habitable Zone}
\label{sec:habitable_zone}

The habitable zone is the annulus around a star within which an
Earth-analog planet --- Earth-mass, Earth-composition, with a
CO$_2$/H$_2$O/N$_2$ atmosphere regulated by a functioning
carbonate--silicate cycle --- can sustain a global body of liquid water
at its surface \parencite{Kasting1993,Kopparapu2013}.  It is bounded
on the inside by the runaway greenhouse limit, at which a
water-saturated atmosphere can no longer radiate the incoming stellar
flux to space, and on the outside by the maximum-greenhouse limit, at
which CO$_2$ condenses out of the atmosphere before the greenhouse
effect can compensate for the reduced insolation.  Both boundaries
are theoretical rather than observed: they emerge from 1D
radiative--convective climate models calibrated to Earth's atmospheric
chemistry.  The concept as it stands has been tested on a single
planetary system --- our own --- and its wider validity is one of the
open questions of astrobiology \parencite{CatlingZahnle2020}.

The outer edge deserves a brief comment because this thesis will
otherwise ignore it.  Unlike the inner edge, the outer boundary does
not represent a catastrophic atmospheric instability but rather the
point at which a self-regulating climate feedback runs out of range.

On timescales longer than ${\sim}\,10^6$ years, atmospheric
CO$_2$ on an Earth-like planet is controlled by the
carbonate--silicate cycle.  Silicate rocks are weathered by
CO$_2$-bearing rainwater; the dissolved products are eventually buried
as carbonate sediments on the ocean floor; subduction and metamorphism
then return CO$_2$ to the atmosphere via volcanism.  The critical
feature of this cycle is that weathering rates are temperature
sensitive: a cooler planet weathers more slowly, so volcanic CO$_2$
accumulates.  A planet receiving less stellar flux therefore responds
by building up a denser CO$_2$ greenhouse, which partly compensates
for the cooling.  This negative feedback acts as a long-term planetary
thermostat \parencite{Kasting1993}.

The thermostat has a physical ceiling.  As the atmospheric CO$_2$
partial pressure rises, two competing effects overtake the warming.
First, the outgoing infrared flux saturates: once the atmosphere is
optically thick at all infrared wavelengths, adding more CO$_2$
cannot trap additional heat.  Second, CO$_2$ is an efficient Rayleigh
scatterer of incoming starlight, so a denser atmosphere reflects a
larger fraction of the incident flux back to space, reducing the
energy available to warm the surface.  The greenhouse warming and
the albedo cooling therefore pull in opposite directions, and the
net heating reaches a maximum --- the maximum-greenhouse limit ---
at a CO$_2$ partial pressure of roughly 8\,bar, beyond which adding
more CO$_2$ cools rather than warms the planet \parencite{Kasting1993}.
For a solar-analog host, \textcite{Kopparapu2013} place this limit at
$S_\mathrm{eff}\approx 0.36$ (${\sim}\,1.67\,\mathrm{AU}$), roughly a
factor of three further from the star than the inner edge.  Once the
thermostat is exhausted the surface freezes, ice raises the planetary
albedo further, and the planet enters a state of runaway glaciation.

Unlike the runaway greenhouse, this transition is gradual rather than
catastrophic.  It can be shifted substantially by non-condensing greenhouse gases.
Primordial H$_2$--He mixtures accreted from a protoplanetary disk can
maintain above-freezing surface temperatures out to 10\,AU around a
G-type star \parencite{PierrehumbertGaidos2011}.  More relevant to
long-lived habitable worlds, \textcite{RamirezKaltenegger2017} showed
that continuous volcanic outgassing of H$_2$ can extend the outer edge
to ${\sim}\,2.4\,\mathrm{AU}$ in our solar system, and CH$_4$ from
volcanism increases the outer-edge distance by ${\sim}\,8\,\%$ for a
solar-analog host (see \textcite{Ramirez2018} for a review).
Most importantly for this thesis, the outer transition leaves no
clean signature on the transit-radius plane.  A cold rocky planet
near the maximum-greenhouse limit has essentially the same bulk radius
as its temperate counterpart: no atmospheric phase transition alters
its envelope mass or structure in a detectable way.  The runaway
greenhouse, by contrast, inflates transit radii by tens of percent
(\Cref{sec:hzied_framework}).  The \HZIED test therefore focuses
exclusively on the inner edge, and the outer boundary appears in
the rest of this thesis only as context.

The inner edge is the focus of this thesis and deserves the more
careful treatment.  Following the seminal calculations of
\textcite{Nakajima1992}, and their revision with modern water-vapour
opacities by \textcite{Goldblatt2013}, the outgoing longwave radiation
of a saturated water-vapour atmosphere asymptotes to a maximum value
as the surface temperature increases --- the Simpson--Nakajima limit
introduced in \Cref{chap:introduction}.  Above this ceiling, no
steady state in radiative--convective equilibrium can balance the
incoming stellar flux while maintaining a liquid ocean: the surface
climbs to silicate-melting temperatures and the planet enters a
steam-atmosphere-over-magma-ocean state.  The 3D radiative-hydrodynamic
simulations of \textcite{Leconte2013}, \textcite{Turbet2021}, and
\textcite{Chaverot2023} refine the numerical value of the threshold
(with clouds, rotation, and day--night asymmetry the threshold shifts
by tens of percent) but confirm the qualitative picture.

Two conversions between the different flux units used in the
literature matter for the rest of the thesis and are worth spelling
out explicitly.  The globally averaged absorbed flux $S_\mathrm{abs}$
is the fraction of top-of-atmosphere stellar radiation not reflected
back to space, averaged over the planet's spherical surface:
%
\begin{equation}
    S_\mathrm{abs} = \frac{1-A}{4}\, S_\mathrm{TOA},
    \label{eq:sabs_stoa}
\end{equation}
%
where $A$ is the Bond albedo and $S_\mathrm{TOA}$ is the
top-of-atmosphere insolation.  The factor of four converts the disk
cross-section that intercepts the stellar beam ($\pi R_p^2$) to the
total surface area of the sphere ($4\pi R_p^2$), and the factor
$(1-A)$ removes the reflected fraction.  For a solar-type host and a
Bond albedo close to Earth's ($A=0.30$), the runaway threshold
$S_\mathrm{thresh}=280\,\mathrm{W\,m^{-2}}$ of absorbed radiation
corresponds to $S_\mathrm{TOA}\approx 1.2\,\Searth$, or equivalently a
semi-major axis $a\approx 0.95\,\mathrm{AU}$.  This matches the
runaway-greenhouse effective flux $S_\mathrm{eff}=1.107$ tabulated by
\textcite{Kopparapu2013} for a solar-analog host once their slightly
different Bond albedo is factored in.

The threshold depends on the spectrum of the incident stellar
radiation.  Cooler M-dwarf hosts emit a larger fraction of their
light at wavelengths at which water vapour absorbs strongly --- the
same coupling that drives the runaway feedback in the first place ---
so the M-dwarf HZ inner edge sits at a lower absorbed flux than the
solar-type edge.  The \textcite{Kopparapu2014} calculations place the
M0V limit at $S_\mathrm{eff}\approx 1.06$ (versus $1.107$ for solar
type), and the shift grows for later M-type hosts.  For the
FGK-dominated \PLATO samples analysed in \Cref{chap:results} the
spectral-type shift is at the few-percent level and can safely be
absorbed into a single fiducial $S_\mathrm{thresh}=280\,\mathrm{W\,m^{-2}}$;
for the M-dwarf-dominated P4 sample it is a larger correction that
motivates the mixed-population variant of \Cref{ssec:hycean_mixed}.

Two qualifications of the concept close the section.  First, the
Earth-analog atmospheric assumption is doing a great deal of the
work.  Variations in planetary rotation rate, cloud fraction,
atmospheric H$_2$/He inventory, or interior degassing history can
shift the HZ boundaries by tens of percent
\parencite{Turbet2021,Chaverot2023}.  Second, and as already
foregrounded in \Cref{chap:introduction}, the entire concept is
calibrated on Earth itself and has never been experimentally tested
as a predictive framework.  The runaway-greenhouse transition, as a
specific prediction of the same climate models, is however
observationally accessible on populations of exoplanets --- and that
is what the rest of this chapter builds up to.


% =======================================================================
\section{The Habitable Zone Inner-Edge Discontinuity}
\label{sec:hzied_framework}

The \HZIED is the statistical signature of the runaway greenhouse
transition on the transit-radius distribution of a small-planet
population.  The claim can be stated in three lines: (i) a wet rocky
planet crossing the runaway threshold $S_\mathrm{thresh}$ inflates
its transit-observable radius by tens of percent (physics laid out in
\Cref{chap:introduction} and \Cref{sec:habitable_zone} above); (ii) at
fixed water inventory this radius change is a step function of
instellation; and (iii) the step survives averaging over an entire
planet population even though individual planet radii are noisy.
Detecting it would constitute the first observational detection of
the runaway greenhouse transition in an exoplanet population.

The signal was first assembled into a testable prediction by
\textcite{Schlecker2024}.  Using the \textsc{Bioverse} pipeline of
\textcite{BixelApai2021}, they built a synthetic transiting-planet
catalogue from a Gaia stellar sample and the \textcite{Bergsten2022}
occurrence-rate priors, injected the radius inflation prescription of
\textcite{Turbet2020} (with the water-retention correction of
\textcite{DornLichtenberg2021}) into hot-side planets with a
probability $f_\mathrm{rgh}$, and asked how large a survey would need
to be to distinguish the resulting step in mean radius from a flat
null.  This thesis inherits that framework and broadens the model
space along two axes: alternative interior-composition prescriptions
\parencite{Aguichine2021} and hydrogen-dominated (``Hycean'')
atmospheres \parencite{Pierrehumbert2022,Innes2023}, both of which
alter the location and amplitude of the step.  The full
implementation --- pipeline, mass--radius grids, water-inventory
priors, and Bayesian test variants --- is deferred to
\Cref{chap:methodology}.

The vocabulary introduced here reappears in every subsequent chapter:
$S_\mathrm{thresh}$ for the runaway threshold, $x_{\mathrm{H_2O}}$
for the water inventory, $f_\mathrm{rgh}$ for the runaway-greenhouse
fraction, and $\ln\Delta Z$ for the log-Bayes factor between the step
hypothesis $H_1$ and the flat null $H_0$.  \Cref{chap:results} plots
$\ln\Delta Z$ against each of the injection parameters in turn.

\Cref{fig:generic_reproduction} closes the section by reproducing
the original \textcite{Schlecker2024} detection on the generic
Gaia-based synthetic sample used by that paper.  It confirms that
the framework recovers the injected radius step at very high
significance under the paper's own assumptions, and provides the
in-principle reference point against which the \PLATO-specific
forecasts of \Cref{chap:results} are compared.

\begin{figure}[H]
    \centering
    \includegraphics[width=0.85\textwidth]{generic_reproduction}
    \caption[Schlecker+2024 reproduction on the generic Gaia sample]{%
        Reproduction of the \textcite{Schlecker2024} baseline
        detection on the generic Gaia-based synthetic sample
        ($N = 500$, $\sigma_R = 2\%$, $\sigma_S = 5\%$).  Injection
        parameters: $S_\mathrm{thresh} = 280\,\mathrm{W\,m^{-2}}$,
        $x_{\mathrm{H_2O}} = 0.005$, $f_\mathrm{rgh} = 0.8$;
        \textcite{Turbet2020} steam grid with the
        \textcite{DornLichtenberg2021} water-retention correction.
        Individual observed radii (crosses), rolling SMA
        ($w = 25$, blue band, empirical per-bin uncertainty), and
        posterior draws from the step model (red curves) are shown.
        The vertical dashed line marks the injected
        $S_\mathrm{thresh}$.  The paper's headline configuration is
        recovered with $\ln\Delta Z \gg 1$, confirming that the
        pipeline reproduces the framework of \textcite{Schlecker2024}
        under identical inputs.}
    \label{fig:generic_reproduction}
\end{figure}

% Figure produced by:
%   Src/figures/generic_reproduction.py
%   → Figures/generated/generic_reproduction.pdf


